\chapter{Suites et séries de fonctions}\label{ch:b2:funcseq}

Lorsque des fonctions convergent vers une fonction, quelles propriétés
survivent au passage à la limite~? La \omterm{def:b2:funcseq:def}{convergence simple} ne préserve
presque rien~; la convergence \emph{uniforme} --- la convergence en \omterm{def:b2:nvs:norm}{norme}
sup --- préserve la continuité, les intégrales sur les segments et, à une
nuance près, les dérivées. Ce chapitre démontre les trois théorèmes de
transfert, leurs versions en séries, et les couronne par le théorème
d'approximation de Weierstrass, démontré grâce aux magnifiques polynômes
probabilistes de Bernstein.

\section{Convergence simple et uniforme}

\begin{definition}\label{def:b2:funcseq:def}
Soient $f_n, f \colon X \to \R$ (ou $\C$, ou un espace normé), $X$ un
ensemble quelconque. $(f_n)$ \omterm{def:b2:integration:improper}{converge} vers $f$
\emph{simplement}\index{convergence simple} lorsque $f_n(x) \to
f(x)$ pour tout $x$~; \emph{uniformément}\index{convergence uniforme}
lorsque
\[
\norm{f_n - f}_\infty = \sup_{x \in X}\, \abs{f_n(x) - f(x)}
\xrightarrow[n \to \infty]{} 0 .
\]
La convergence uniforme entraîne la convergence simple~; sur
$C(\intcc{a}{b})$, la convergence uniforme est exactement la convergence
dans l'\omterm{def:b2:nvs:banach}{espace de Banach}
$\bigl(C(\intcc{a}{b}), \norm\cdot_\infty\bigr)$ du
\cref{ch:b2:nvs}.
\end{definition}

\begin{example}\label{ex:b2:funcseq:xn}
Sur $\intcc{0}{1}$, $f_n(x) = x^n$ \omterm{def:b2:integration:improper}{converge} \omterm{def:b2:funcseq:def}{simplement} vers la limite
\emph{discontinue} $f = \mathbf{1}_{\{1\}}$~; la convergence
n'est pas uniforme~: $\norm{f_n - f}_\infty \geq f_n\bigl(1 -
\tfrac1n\bigr) = (1 - \tfrac1n)^n \to \eu^{-1} \neq 0$. Sur
$\intcc{0}{a}$ avec $a < 1$ elle l'\emph{est} ($\sup = a^n \to
0$)~: l'uniformité est une propriété du domaine autant que de la
suite.
\end{example}

\begin{omfigure}
\begin{tikzpicture}
\begin{axis}[omaxis, width=8.2cm, height=5.6cm,
    xmin=0, xmax=1.08, ymin=0, ymax=1.15,
    xtick={0,1}, ytick={1}]
  \addplot[omDef!40, thick, domain=0:1, samples=80] {x^2};
  \addplot[omDef!60, thick, domain=0:1, samples=80] {x^5};
  \addplot[omDef, thick, domain=0:1, samples=120] {x^15};
  \addplot[omThm, very thick, domain=0:1, samples=160] {x^40};
  \node[font=\small, omDef] at (axis cs:0.62,0.5) {$x^2$};
  \node[font=\small, omDef] at (axis cs:0.83,0.48) {$x^5$};
  \node[font=\small, omThm] at (axis cs:0.987,0.5) {$x^{40}$};
\end{axis}
\end{tikzpicture}

{\small La suite $x^n$ sur $\intcc{0}{1}$~: les graphes s'affaissent
vers $0$ mais doivent tous grimper jusqu'à $1$ en $x = 1$ --- la
distance sup à la limite simple discontinue ne descend jamais sous une
constante.}
\end{omfigure}

\begin{example}[Deux limites qui refusent de commuter]\label{ex:b2:funcseq:doublelimit}
Tout le chapitre porte sur l'interversion de limites, voici donc le
plus petit échec possible. Posons $a_{n,m} = \dfrac{n}{n+m}$ pour
$n, m \geq 1$. Alors
\[
\lim_{m\to\infty}\Bigl(\lim_{n\to\infty}a_{n,m}\Bigr)
= \lim_{m\to\infty} 1 = 1,
\qquad
\lim_{n\to\infty}\Bigl(\lim_{m\to\infty}a_{n,m}\Bigr)
= \lim_{n\to\infty} 0 = 0 :
\]
les deux limites itérées existent et diffèrent. Chaque théorème de
transfert de ce chapitre est un permis d'intervertir deux limites ---
$\lim_n$ avec $\lim_{x\to a}$ (continuité), avec $\int$
(intégration), avec $\frac{\dd}{\dd x}$ (dérivation) --- et la
\omterm{def:b2:funcseq:def}{convergence uniforme} est précisément le prix qui rend l'interversion
licite. Éclairage final~: chaque fois qu'une «~démonstration~»
échange silencieusement deux opérations de limite, ce tableau de deux
lignes est le contre-exemple à lui opposer~; les bosses glissantes de
l'\cref{exo:b2:funcseq:2} sont le même phénomène affublé d'un signe
intégral.
\end{example}

\section{Les trois théorèmes de transfert}

\begin{theorem}[Continuité]\label{thm:b2:funcseq:continuity}
Si chaque $f_n$ est \omterm{def:b2:metric:continuity}{continue} en $a$ et $f_n \to f$ \omterm{def:b2:funcseq:def}{uniformément} sur un
voisinage de $a$, alors $f$ est \omterm{def:b2:metric:continuity}{continue} en $a$. Une limite uniforme
de fonctions \omterm{def:b2:metric:continuity}{continues} est \omterm{def:b2:metric:continuity}{continue}.
\end{theorem}

\begin{proof}
L'argument en $3\varepsilon$ déjà utilisé dans le
\cref{thm:b2:metric:rncomplete}~: on choisit $n$ tel que $\norm{f_n -
f}_\infty \leq \varepsilon$, puis $\delta$ par la continuité de
$f_n$ en $a$~; pour $\abs{x - a} \leq \delta$,
\[
\abs{f(x) - f(a)} \leq \abs{f(x) - f_n(x)} + \abs{f_n(x) - f_n(a)}
+ \abs{f_n(a) - f(a)} \leq 3\varepsilon . \qedhere
\]
\end{proof}

\begin{example}[L'uniformité échoue exactement là où la limite se brise]\label{ex:b2:funcseq:threshold}
Sur $\intcc{0}{2}$, posons $f_n(x) = \dfrac{x^n}{1 + x^n}$. La
limite simple est une fonction en trois morceaux~:
\[
f(x) = \begin{cases} 0 & 0 \leq x < 1,\\[2pt]
\tfrac12 & x = 1,\\[2pt]
1 & 1 < x \leq 2, \end{cases}
\]
discontinue en $1$, donc d'après le \cref{thm:b2:funcseq:continuity}
la convergence ne peut être uniforme sur $\intcc{0}{2}$. Sur les
morceaux fermés évitant le seuil elle l'est~: pour $0 \leq x \leq
a < 1$,
\[
\sup_{\intcc{0}{a}}\abs{f_n - 0}
= \frac{a^n}{1 + a^n} \leq a^n \to 0 ,
\]
et pour $1 < b \leq x \leq 2$,
\[
\sup_{\intcc{b}{2}}\abs{f_n - 1}
= \frac{1}{1 + b^n} \leq b^{-n} \to 0 ,
\]
les deux suprema étant calculés par monotonie de $u \mapsto
\frac{u}{1+u}$ et de $x \mapsto x^n$. Éclairage final~: l'échec de
l'uniformité est localisé à la discontinuité de la limite --- la même
géométrie que l'\cref{ex:b2:funcseq:xn}, et la raison pour laquelle la
discipline «~uniforme sur tout segment intérieur~» revient tout au
long du chapitre.
\end{example}

\begin{theorem}[Intégration sur un segment]\label{thm:b2:funcseq:integration}
Si $f_n \to f$ \omterm{def:b2:funcseq:def}{uniformément} sur $\intcc{a}{b}$, avec $f_n$ \omterm{def:b2:metric:continuity}{continue}
par morceaux ($f$ de même), alors
\[
\int_a^b f_n \longrightarrow \int_a^b f .
\]
\end{theorem}

\begin{proof}
La linéarité et l'inégalité triangulaire pour les intégrales donnent
\[
\Bigl|\int_a^b f_n - \int_a^b f\Bigr|
= \Bigl|\int_a^b (f_n - f)\Bigr|
\leq \int_a^b\abs{f_n - f}
\leq (b - a)\,\norm{f_n - f}_\infty \longrightarrow 0 .
\]
Le facteur de longueur $(b - a)$ est l'endroit où la compacité du
segment intervient~: sur les intervalles non \omterm{def:b2:metric:compact}{compacts}, la même
estimation produit la borne inutile $\infty\cdot0$, et la conclusion
échoue véritablement sans domination --- les bosses plates $f_n =
\frac1n\mathbf 1_{\intcc{0}{n}}$ convergent \omterm{def:b2:funcseq:def}{uniformément} vers $0$ sur
$\intco{0}{\infty}$ tout en conservant $\int f_n = 1$ (remarque sur
les pièges ci-dessous), et les bosses glissantes de la remarque du
\cref{ch:b2:integration} font de même avec la \omterm{def:b2:funcseq:def}{convergence simple}~; la
\omterm{def:b2:funcseq:def}{convergence uniforme} contrôle les hauteurs, jamais les largeurs.
\end{proof}

\begin{omfigure}
\begin{tikzpicture}
\begin{axis}[omaxis, width=8.6cm, height=5.4cm,
    xmin=0, xmax=1.05, ymin=0, ymax=1.9,
    xtick={0,1}, ytick={1}]
  \addplot[omDef!45, thick, domain=0:1.05, samples=120]
    {2*x*exp(-2*x^2)};
  \addplot[omDef!70, thick, domain=0:1.05, samples=120]
    {5*x*exp(-5*x^2)};
  \addplot[omDef, thick, domain=0:1.05, samples=160]
    {15*x*exp(-15*x^2)};
  \node[font=\small, omDef] at (axis cs:0.62,0.42) {$n=2$};
  \node[font=\small, omDef] at (axis cs:0.40,0.80) {$n=5$};
  \node[font=\small, omDef] at (axis cs:0.24,1.55) {$n=15$};
\end{axis}
\end{tikzpicture}

{\small Les bosses $g_n(x) = nx\,\eu^{-nx^2}$ de
l'\cref{exo:b2:funcseq:1}~: elles convergent vers $0$ en tout point,
mais les pics (de hauteur $\sim\sqrt{n/(2\eu)}$, dérivant vers
$0$) croissent sans borne --- \omterm{def:b2:funcseq:def}{convergence simple} avec
$\norm{g_n}_\infty \to \infty$, et $\int_0^1 g_n \to \frac12
\neq 0$~: la masse se cache sous le pic mobile.}
\end{omfigure}

\begin{theorem}[Dérivation]\label{thm:b2:funcseq:differentiation}
Soit $f_n$ de classe $C^1$ sur un intervalle $I$, avec~: $(f_n')$
convergeant \emph{\omterm{def:b2:funcseq:def}{uniformément}} sur $I$ (ou sur tout segment de $I$)
vers un $g$, et $(f_n(x_0))$ convergeant en un point $x_0$. Alors
$(f_n)$ \omterm{def:b2:integration:improper}{converge} (\omterm{def:b2:funcseq:def}{uniformément} sur les segments) vers une fonction $f$
de classe $C^1$, et $f' = g$~: on peut dériver la limite.
\end{theorem}

\begin{proof}
Définissons $f(x) = \lim f_n(x_0) + \int_{x_0}^x g$~: légitime, $g$
étant \omterm{def:b2:metric:continuity}{continue} --- en effet $g$ est la limite \emph{uniforme} sur les
segments des $f_n'$ \omterm{def:b2:metric:continuity}{continues}, donc le
\cref{thm:b2:funcseq:continuity} s'applique, et l'intégrale d'une
fonction \omterm{def:b2:metric:continuity}{continue} est bien définie avec, par le théorème fondamental
de l'analyse,
\[
f'(x) = g(x) \qquad (x \in I) :
\]
la limite candidate est de classe $C^1$ avec la bonne dérivée
\emph{par construction}, avant toute preuve de convergence. Par le
théorème fondamental de nouveau, $f_n(x) = f_n(x_0) + \int_{x_0}^x
f_n'$~; en soustrayant,
\[
\abs{f_n(x) - f(x)} \leq \abs{f_n(x_0) - \lim f_n(x_0)}
+ \abs{x - x_0}\,\norm{f_n' - g}_{\infty} ,
\]
qui tend vers $0$ \omterm{def:b2:funcseq:def}{uniformément} sur tout segment. Et $f$ est de classe
$C^1$ avec $f' = g$ par construction.
\end{proof}

\begin{example}[Pourquoi l'hypothèse porte sur les dérivées]\label{ex:b2:funcseq:sqrtabs}
Posons $F_n(x) = \sqrt{x^2 + \frac1n}$ sur $\R$. Chaque $F_n$ est
de classe $C^1$ (et même $C^\infty$), et la convergence vers $\abs x$
est uniforme sur tout $\R$~:
\[
0 \leq F_n(x) - \abs x
= \frac{(x^2 + \frac1n) - x^2}{\sqrt{x^2+\frac1n} + \abs x}
= \frac{1/n}{\sqrt{x^2 + \frac1n} + \abs x}
\leq \frac{1/n}{1/\sqrt n} = \frac{1}{\sqrt n} .
\]
Pourtant la limite $\abs x$ n'est pas dérivable en $0$~: la
\omterm{def:b2:funcseq:def}{convergence uniforme} des \emph{fonctions}, si rapide soit-elle, ne
transfère aucune dérivabilité. L'échec est visible sur les dérivées~:
\[
F_n'(x) = \frac{x}{\sqrt{x^2 + \frac1n}}
\longrightarrow \begin{cases} 1 & x > 0,\\ 0 & x = 0,\\
-1 & x < 0, \end{cases}
\]
une limite simple discontinue, donc $(F_n')$ ne peut converger
\omterm{def:b2:funcseq:def}{uniformément} près de $0$ (le \cref{thm:b2:funcseq:continuity} de
nouveau). Éclairage final~:
le \cref{thm:b2:funcseq:differentiation} suppose délibérément la
\omterm{def:b2:funcseq:def}{convergence uniforme} des $f_n'$, non des $f_n$ ---
cet exemple en est la raison.
\end{example}

\section{Séries de fonctions}

\begin{definition}\label{def:b2:funcseq:series}
Une série de fonctions $\sum u_n$ \omterm{def:b2:integration:improper}{converge} \omterm{def:b2:funcseq:def}{simplement}/\omterm{def:b2:funcseq:def}{uniformément}
lorsque ses sommes partielles le font. Elle \omterm{def:b2:integration:improper}{converge}
\emph{normalement}\index{convergence normale} (sur $X$) lorsque $\sum
\norm{u_n}_\infty < \infty$. La convergence normale entraîne la
\omterm{def:b2:funcseq:def}{convergence uniforme} (dans l'\omterm{def:b2:nvs:banach}{espace de Banach} des fonctions bornées~:
\cref{thm:b2:nvs:absoluteconvergence}), qui entraîne la \omterm{def:b2:funcseq:def}{convergence
simple}~; les deux implications sont strictes.
\end{definition}

\begin{example}[Une série, trois verdicts]\label{ex:b2:funcseq:threemodes}
Prenons $u_n(x) = \frac{x^n}{n}$ sur $\intco{0}{1}$. \emph{\omterm{def:b2:funcseq:def}{Simplement}~:}
\omterm{def:b2:integration:improper}{converge} pour tout $x \in \intco01$ (comparaison avec la
série géométrique). \emph{\omterm{def:b2:funcseq:series}{Normalement} sur $\intcc{0}{a}$, $a < 1$~:}
$\norm{u_n}_{\infty,\intcc0a} = \frac{a^n}{n}$, \omterm{def:b2:series:summable}{sommable}.
\emph{Pas \omterm{def:b2:funcseq:series}{normalement} sur $\intco{0}{1}$~:}
$\norm{u_n}_{\infty,\intco01} = \frac1n$, et $\sum\frac1n$
diverge. \emph{Pas même \omterm{def:b2:funcseq:def}{uniformément} sur $\intco{0}{1}$~:} le
reste résiste près de $1$,
\[
R_N(x) = \sum_{n>N}\frac{x^n}{n}
\geq \sum_{n=N+1}^{2N}\frac{x^n}{n}
\geq \frac{N\,x^{2N}}{2N} = \frac{x^{2N}}{2}
\xrightarrow[x\to1^-]{} \frac12 ,
\]
donc $\sup_{\intco01}\abs{R_N} \geq \frac12$ pour tout $N$.
Éclairage final~: les quatre verdicts coexistent paisiblement --- la
somme $-\ln(1-x)$ est \omterm{def:b2:metric:continuity}{continue} sur $\intco{0}{1}$ car la
continuité ne requiert l'uniformité que \emph{près de chaque point},
c.-à-d.\ sur les segments $\intcc0a$~; exploser au bord est le droit
de la somme.
\end{example}

\begin{theorem}[Transfert pour les séries]\label{thm:b2:funcseq:seriestransfer}
Si $\sum u_n$ \omterm{def:b2:integration:improper}{converge} \omterm{def:b2:funcseq:def}{uniformément} (par exemple \omterm{def:b2:funcseq:series}{normalement}) sur
l'ensemble concerné~: la continuité de tous les $u_n$ en $a$ passe à
la somme~; l'intégration sur un segment peut se faire terme à terme~;
et si $\sum u_n(x_0)$ \omterm{def:b2:integration:improper}{converge} tandis que $\sum u_n'$ \omterm{def:b2:integration:improper}{converge}
\omterm{def:b2:funcseq:def}{uniformément} sur les segments, la somme est de classe $C^1$ de dérivée
$\sum u_n'$.
\end{theorem}

\begin{proof}
Tout est le théorème correspondant appliqué aux sommes partielles
$S_N = \sum_{n\leq N}u_n$, qui sont des sommes finies de fonctions
ayant la régularité voulue. Continuité~: chaque $S_N$ est \omterm{def:b2:metric:continuity}{continue} en
$a$ et $S_N \to \sum u_n$ \omterm{def:b2:funcseq:def}{uniformément}~:
\cref{thm:b2:funcseq:continuity}. Intégration~: sur le segment,
\[
\int_a^b \sum_{n\geq0} u_n
= \lim_N \int_a^b S_N
= \lim_N \sum_{n=0}^{N}\int_a^b u_n
= \sum_{n\geq0}\int_a^b u_n
\]
par le \cref{thm:b2:funcseq:integration} (première égalité) et la
linéarité de l'intégrale (seconde). Dérivation~: les
$S_N$ sont de classe $C^1$, $S_N(x_0)$ \omterm{def:b2:integration:improper}{converge}, et $S_N' =
\sum_{n\leq N}u_n'$ \omterm{def:b2:integration:improper}{converge} \omterm{def:b2:funcseq:def}{uniformément} sur les segments~: le
\cref{thm:b2:funcseq:differentiation} donne que la somme est de
classe $C^1$ de dérivée $\lim S_N' = \sum u_n'$.
\end{proof}

\begin{remark}[Pièges classiques]
Quatre pièges, tous vus dans des copies d'examen. \emph{(i)
Suprema à moitié vérifiés~:} évaluer $f_n$ le long d'une suite
$x_n$ bien choisie ne fournit qu'une \emph{minoration} de
$\norm{f_n - f}_\infty$ --- suffisante pour infirmer l'uniformité
(comme dans l'\cref{ex:b2:funcseq:xn}), jamais pour l'établir~; pour
l'établir, majorer le sup par un calcul valable pour \emph{tout} $x$.
\emph{(ii) Uniformité sur le mauvais ensemble~:} la \omterm{def:b2:funcseq:series}{convergence
normale} ou uniforme vaut souvent sur tout $\intcc{-a}{a}$ ou
$\intco\delta\infty$ mais échoue sur la réunion ouverte~; ce n'est pas
un obstacle --- la continuité et la dérivabilité sont locales, donc la
discipline segment par segment de l'\cref{ex:b2:funcseq:zeta} les
donne sur tout l'\omterm{def:b2:metric:topology}{ouvert}. \emph{(iii) Intégrer sur des non-segments~:}
le \cref{thm:b2:funcseq:integration} est un énoncé sur les segments~;
sur $\intco0\infty$, la \omterm{def:b2:funcseq:def}{convergence uniforme} n'empêche pas la masse de
s'échapper à l'infini ($f_n =
\frac1n\mathbf 1_{\intcc{0}{n}}$ \omterm{def:b2:integration:improper}{converge} \omterm{def:b2:funcseq:def}{uniformément} vers $0$,
avec $\int f_n = 1$) --- utiliser là la convergence dominée.
\emph{(iv) Dériver la limite~:}
\cref{ex:b2:funcseq:sqrtabs}~; l'hypothèse de dérivée porte sur
$(f_n')$, et aucune vitesse de convergence de $(f_n)$ ne peut la
remplacer.
\end{remark}

\begin{example}[La fonction $\zeta$ de Riemann]\label{ex:b2:funcseq:zeta}
$\zeta(s) = \sum_{n\geq1} n^{-s}$ \omterm{def:b2:integration:improper}{converge} \omterm{def:b2:funcseq:series}{normalement} sur toute
demi-droite $\intco{a}{+\infty}$, $a > 1$ ($\norm{n^{-s}}_\infty =
n^{-a}$, \omterm{def:b2:series:summable}{sommable})~: $\zeta$ est \omterm{def:b2:metric:continuity}{continue} sur $\intoo{1}{+\infty}$~;
en dérivant terme à terme (la série dérivée $\sum -\ln n\; n^{-s}$
\omterm{def:b2:integration:improper}{converge} elle aussi \omterm{def:b2:funcseq:series}{normalement} sur $\intco{a}{\infty}$), $\zeta$ est
de classe $C^1$ --- et, par itération, $C^\infty$ --- avec
$\zeta'(s) = -\sum \frac{\ln n}{n^s}$. Noter la discipline~: la
\omterm{def:b2:funcseq:series}{convergence normale} se vérifie sur des \emph{sous}-demi-droites,
jamais sur l'\omterm{def:b2:metric:topology}{ouvert} $\intoo{1}{\infty}$ lui-même, où elle échoue.
\end{example}

\begin{example}[Une série logarithmique, traitée jusqu'au bout]\label{ex:b2:funcseq:logseries}
Posons $F(x) = \sum_{n\geq1} \frac{\eu^{-nx}}{n}$ sur
$\intoo{0}{\infty}$. Chaque terme est borné sur $\intco{\delta}
\infty$ par $\frac{\eu^{-n\delta}}{n} \leq \eu^{-n\delta}$, une
série géométrique convergente~: \omterm{def:b2:funcseq:series}{convergence normale} sur tout
$\intco\delta\infty$, donc $F$ est \omterm{def:b2:metric:continuity}{continue} sur
$\intoo{0}{\infty}$. La série dérivée $\sum -\eu^{-nx}$ est
de même \omterm{def:b2:funcseq:series}{normalement} convergente sur $\intco\delta\infty$
($\norm{\eu^{-nx}}_{\infty,\intco\delta\infty} =
\eu^{-n\delta}$), donc $F$ est de classe $C^1$ de dérivée géométrique~:
\[
F'(x) = -\sum_{n\geq1}\eu^{-nx}
= \frac{-\eu^{-x}}{1 - \eu^{-x}}
= \frac{-1}{\eu^{x} - 1} .
\]
Par itération, $F$ est de classe $C^\infty$. En intégrant $F'$ (aussi
bien $F$ que $x \mapsto -\ln(1 - \eu^{-x})$ s'annulent en $+\infty$ et
ont la même dérivée sur $\intoo0\infty$)~:
\[
F(x) = -\ln\bigl(1 - \eu^{-x}\bigr),
\]
la série logarithmique en $t = \eu^{-x}$. Éclairage final~: lorsque $x
\to 0^+$, $F(x) = -\ln(x + O(x^2)) = \ln\frac1x + O(x)$ --- la
série diverge logarithmiquement au bord, exactement comme la série
harmonique qu'elle devient en $x = 0$~; la \omterm{def:b2:funcseq:series}{convergence normale} sur
$\intco\delta\infty$ mais non sur $\intoo0\infty$ en est le
symptôme.
\end{example}

\begin{method}[Prouver ou réfuter la convergence uniforme]\label{met:b2:funcseq:uniform}
Pour $f_n \to f$ \omterm{def:b2:funcseq:def}{simplement} sur $X$~:
\begin{enumerate}
  \item Calculer ou majorer $\norm{f_n - f}_\infty$~: étudier la
        fonction $x \mapsto \abs{f_n(x) - f(x)}$ (dérivée,
        monotonie) pour localiser son maximum~; une majoration
        valable pour tout $x$ qui tend vers $0$ prouve l'uniformité.
  \item Pour \emph{réfuter}~: exhiber des points $x_n$ avec
        $\abs{f_n(x_n) - f(x_n)} \not\to 0$ (souvent $x_n$ suit
        la bosse mobile, comme dans l'\cref{exo:b2:funcseq:1})~; ou
        invoquer un théorème de transfert par contraposée --- une
        limite discontinue de fonctions \omterm{def:b2:metric:continuity}{continues}
        (\cref{ex:b2:funcseq:threshold}), ou $\int f_n
        \not\to \int f$ sur un segment.
  \item Pour les séries, tenter d'abord la \omterm{def:b2:funcseq:series}{convergence normale}
        ($\sum\sup\abs{u_n} < \infty$)~; si elle échoue globalement,
        la tester sur les sous-segments qui comptent
        (\cref{ex:b2:funcseq:threemodes})~; si elle échoue
        partout, la \omterm{def:b2:funcseq:def}{convergence uniforme} peut encore valoir via la
        majoration du reste alterné
        (\cref{exo:b2:funcseq:4}) ou la sommation par parties.
\end{enumerate}
\end{method}

\section{Le théorème d'approximation de Weierstrass}

\begin{theorem}[Weierstrass, via Bernstein]\label{thm:b2:funcseq:weierstrass}
Toute fonction \omterm{def:b2:metric:continuity}{continue} $f \colon \intcc{0}{1} \to \R$ est limite
uniforme de polynômes --- explicitement, de ses \emph{polynômes de
Bernstein}\index{polynômes de Bernstein}\index{théorème d'approximation de Weierstrass}
\[
B_n(f)(x) = \sum_{k=0}^{n} f\Bigl(\frac kn\Bigr)\binom nk x^k
(1-x)^{n-k} .
\]
\end{theorem}

\begin{proof}
Fixons $x \in \intcc{0}{1}$ et posons $p_k(x) = \binom nk x^k(1 -
x)^{n-k}$. Trois identités binomiales, obtenues en évaluant
$(x + y)^n$ et ses deux dérivées en $x$ au point $y = 1 - x$~:
\[
\sum_k p_k = 1,
\qquad
\sum_k k\,p_k = nx,
\qquad
\sum_k k(k-1) p_k = n(n-1)x^2 .
\]
En détail~: $(x+y)^n = \sum_k\binom nk x^ky^{n-k}$ en $y = 1-x$
est la première~; en dérivant en $x$,
\[
n(x+y)^{n-1} = \sum_k k\binom nk x^{k-1}y^{n-k} ,
\]
puis en multipliant par $x$ et en posant $y = 1 - x$ on obtient la
deuxième~; en dérivant deux fois et en multipliant par $x^2$ on
obtient la troisième. En développant $(k - nx)^2 = k(k-1) + k(1 - 2nx) +
n^2x^2$ et en combinant les trois~:
\[
\sum_k (k - nx)^2 p_k
= n(n-1)x^2 + nx(1 - 2nx) + n^2x^2
= nx(1 - x) \leq \frac n4 ,
\]
l'\emph{identité de variance}.

Estimons maintenant, en utilisant $\sum p_k = 1$~:
\[
\abs{B_n(f)(x) - f(x)}
\leq \sum_{k} \Bigl| f\Bigl(\frac kn\Bigr) - f(x)\Bigr|\, p_k(x)
= \Sigma_{\text{proche}} + \Sigma_{\text{loin}} ,
\]
en séparant selon que $\abs{\frac kn - x} \leq \delta$ ou non.
Étant donné $\varepsilon > 0$, la continuité uniforme de $f$ (Heine)
fournit $\delta$ avec $\Sigma_{\text{proche}} \leq \varepsilon$. Pour la
somme lointaine, avec $M = \norm f_\infty$~: par l'identité de variance
et l'astuce de comptage de Tchebychev,
\[
\Sigma_{\text{loin}} \leq 2M \sum_{\abs{k - nx} > n\delta} p_k
\leq 2M\,\frac{\sum_k (k - nx)^2 p_k}{n^2\delta^2}
\leq \frac{2M}{4 n \delta^2} \cdot 1
= \frac{M}{2n\delta^2}
\xrightarrow[n\to\infty]{} 0 ,
\]
\omterm{def:b2:funcseq:def}{uniformément} en $x$. Donc $\norm{B_n(f) - f}_\infty \leq \varepsilon +
\frac{M}{2n\delta^2} \leq 2\varepsilon$ pour $n$ grand.
\end{proof}

\begin{remark}
Par substitution affine le théorème vaut sur tout segment
$\intcc{a}{b}$. Il échoue sur $\R$ (une limite uniforme de polynômes
sur $\R$ est un polynôme~: \cref{exo:b2:funcseq:8}). La lecture
probabiliste --- $B_n(f)(x)$ est l'espérance de $f$ en une moyenne
binomiale, et la borne de variance est l'inégalité de Tchebychev ---
est rendue honnête au \cref{ch:b2:genfun}.
\end{remark}

\begin{omfigure}
\begin{tikzpicture}
\begin{axis}[omaxis, width=8.6cm, height=6.0cm,
    xmin=0, xmax=1.05, ymin=0, ymax=1.08,
    xtick={0,1}, ytick={1}]
  \addplot[omThm, very thick, domain=0:1, samples=80] {x^2};
  \addplot[omDef!45, thick, domain=0:1, samples=2] {x};
  \addplot[omDef!70, thick, domain=0:1, samples=80]
    {x^2 + x*(1-x)/2};
  \addplot[omDef, thick, domain=0:1, samples=80]
    {x^2 + x*(1-x)/6};
  \node[font=\small, omDef] at (axis cs:0.30,0.37) {$B_1f$};
  \node[font=\small, omDef] at (axis cs:0.52,0.46) {$B_2f$};
  \node[font=\small, omDef] at (axis cs:0.76,0.56) {$B_6f$};
  \node[font=\small, omThm] at (axis cs:0.88,0.68) {$f$};
\end{axis}
\end{tikzpicture}

{\small Approximation de Bernstein de $f(x) = x^2$ (en rouge),
utilisant la formule exacte $B_nf = x^2 + \frac{x(1-x)}{n}$ de
l'\cref{exo:b2:funcseq:7}~: $B_1f$ est la corde, et chaque doublement
de $n$ divise l'écart par deux. Fiable mais lent --- la saturation en
$\frac1n$ que le théorème de Voronovskaya (problème du week-end) rend
exacte.}
\end{omfigure}

\begin{remark}[Où cela sert]
L'approximation de Weierstrass est le théorème de densité de l'analyse
classique~: elle rend $C(\intcc ab)$ séparable, permet de vérifier des
identités intégrales sur les seuls polynômes (problèmes de moments),
et sous-tend la version trigonométrique démontrée dans le chapitre de
Fourier via le noyau de Fejér. Le problème du week-end de ce chapitre
extrait le contenu quantitatif de la preuve de Bernstein --- des
vitesses de convergence régies par le module de continuité --- puis
isole ce qui l'a vraiment fait fonctionner, dans le théorème de
Korovkin~: positivité et trois fonctions test. Le volume de troisième
année généralise l'énoncé de densité à des \omterm{def:b2:structures:algebra}{sous-algèbres} quelconques
(Stone--Weierstrass) et aux espaces \omterm{def:b2:metric:compact}{compacts}.
\end{remark}

\begin{example}[Approximation polygonale, avec une vitesse]\label{ex:b2:funcseq:polygonal}
Pour $f$ $L$-lipschitzienne sur $\intcc{0}{1}$, soit $I_nf$
l'interpolant affine par morceaux aux nœuds $\frac kn$. Sur une
maille $\intcc{\frac kn}{\frac{k+1}n}$, à la fois $f(x)$ et
$I_nf(x)$ se situent entre les valeurs extrêmes qu'une fonction
$L$-lipschitzienne peut prendre étant données les deux valeurs
nodales, donc pour $x$ dans la maille, en posant $x_k = \frac kn$~:
\[
\abs{I_nf(x) - f(x)}
\leq \abs{I_nf(x) - f(x_k)} + \abs{f(x_k) - f(x)}
\leq L\,\abs{x - x_k} + L\,\abs{x - x_k}
\leq \frac{2L}{n}
\]
(l'interpolant est lui-même $L$-lipschitzien sur la maille~: sa
pente est un taux d'accroissement de $f$). D'où $\norm{I_nf -
f}_\infty \leq \frac{2L}{n}$~: l'approximation polygonale des
fonctions \omterm{def:b2:metric:continuity}{lipschitziennes} \omterm{def:b2:integration:improper}{converge} à la vitesse $\frac1n$ ---
\emph{plus rapide} que le $\frac{1}{\sqrt n}$ de Bernstein pour la
même classe (problème du week-end, Partie~II). Éclairage final~: le
polygone interpole mais n'est pas régulier, Bernstein est régulier
mais lent~; il n'y a pas de repas gratuit entre régularité de
l'approximant et vitesse --- un compromis rendu précis par les
résultats de saturation du problème du week-end.
\end{example}

\begin{remark}[Perspectives dans ce volume]
La \omterm{def:b2:funcseq:def}{convergence uniforme} est le cheval de bataille de ce livre à partir
d'ici. Le chapitre sur les séries entières fonctionne entièrement sur
la \omterm{def:b2:funcseq:series}{convergence normale} sur les sous-disques \omterm{def:b2:metric:compact}{compacts} --- chaque
théorème terme à terme qui y figure est un cas particulier des
théorèmes de transfert de ce chapitre. Le chapitre de Fourier vit un
étage au-dessus~: ses sommes partielles $S_N$ échouent exactement là
où ce chapitre prévient qu'elles le pourraient (simple mais non
uniforme aux sauts), et ses moyennes de Fejér réussissent par la même
mécanique en $3\varepsilon$ qui a prouvé le
\cref{thm:b2:funcseq:continuity}. Le chapitre sur les équations
différentielles définit $\eu^{tA}$ par une série \omterm{def:b2:funcseq:series}{normalement}
convergente et la dérive terme à terme --- littéralement le
\cref{thm:b2:funcseq:seriestransfer} appliqué aux coefficients
matriciels. En cas de doute plus loin dans le livre sur «~pourquoi
peut-on faire cela~», la réponse est en général un théorème de ce
chapitre.
\end{remark}

\section{Exercices}

\begin{exercise}[$\star$]\label{exo:b2:funcseq:1}
Étudier la \omterm{def:b2:funcseq:def}{convergence simple} et uniforme sur $\intcc{0}{1}$, puis
sur $\intcc{0}{a}$ ($a < 1$) ou $\intco{\delta}{1}$ selon le cas, de~:
\[
f_n(x) = \frac{x}{1 + nx},
\qquad
g_n(x) = n x\,\eu^{-n x^2},
\qquad
h_n(x) = x^n(1 - x^n).
\]
\end{exercise}

\begin{exercise}[$\star$]\label{exo:b2:funcseq:2}
Prouver que $\displaystyle\int_0^1 g_n \not\to \int_0^1 \lim g_n$
pour les $g_n$ de l'\cref{exo:b2:funcseq:1}, et réconcilier avec le
\cref{thm:b2:funcseq:integration}.
\end{exercise}

\begin{exercise}[$\star$]\label{exo:b2:funcseq:3}
Prouver que $S(x) = \sum_{n\geq1} \dfrac{x^n}{n^2}$ est \omterm{def:b2:metric:continuity}{continue} sur
$\intcc{-1}{1}$, et que $S$ est de classe $C^1$ sur $\intoo{-1}{1}$
avec $S'(x) = -\frac{\ln(1-x)}{x}$ pour $0 < \abs x < 1$.
\end{exercise}

\begin{exercise}[$\star\star$]\label{exo:b2:funcseq:4}
Soit $F(x) = \sum_{n \geq 0} \dfrac{(-1)^n}{n + x}$ sur
$\intoo{0}{+\infty}$. Prouver la \omterm{def:b2:funcseq:def}{convergence uniforme} (non normale)
sur $\intco{\delta}{\infty}$ via la majoration du reste des séries
\omterm{def:b2:linalg:alternating}{alternées}, la continuité, et l'équation fonctionnelle $F(x) + F(x + 1) =
\frac1x$.
\end{exercise}

\begin{exercise}[$\star\star$]\label{exo:b2:funcseq:5}
(Dini) Soit $f_n \colon K \to \R$ \omterm{def:b2:metric:continuity}{continue} sur un \omterm{def:b2:metric:def}{espace métrique}
\emph{\omterm{def:b2:metric:compact}{compact}}, avec $f_n \to f$ \emph{\omterm{def:b2:funcseq:def}{simplement}}, $f$ \omterm{def:b2:metric:continuity}{continue},
et $(f_n(x))$ \emph{décroissante} en $n$ pour chaque $x$. Prouver que
la convergence est uniforme. \emph{(Étant donné $\varepsilon$, les
\omterm{def:b2:metric:topology}{ouverts} $U_n = \{x : f_n(x) - f(x) < \varepsilon\}$ croissent et
recouvrent $K$~; extraire un sous-recouvrement fini ---
\cref{thm:b2:metric:borellebesgue}.)}
\end{exercise}

\begin{exercise}[$\star\star$]\label{exo:b2:funcseq:6}
Prouver que $\displaystyle\lim_{n\to\infty} \int_0^1
\frac{n\,f(x)}{1 + n^2x^2}\,\dd x = \frac{\pi}{2} f(0)$ pour toute
fonction $f$ \omterm{def:b2:metric:continuity}{continue} sur $\intcc{0}{1}$. \emph{(Substituer $u = nx$~;
isoler $f(0)$~; dominer.)}
\end{exercise}

\begin{exercise}[$\star\star$]\label{exo:b2:funcseq:7}
Calculer explicitement les \omterm{thm:b2:funcseq:weierstrass}{polynômes de Bernstein} de $f(x) = x^2$ et
vérifier l'erreur uniforme $\norm{B_n f - f}_\infty =
O\bigl(\frac1n\bigr)$ prédite par la preuve du
\cref{thm:b2:funcseq:weierstrass} --- ici exactement $\frac{x(1 -
x)}{n}$ en chaque point.
\end{exercise}

\begin{exercise}[$\star\star$]\label{exo:b2:funcseq:8}
Prouver que si des polynômes $P_n$ convergent \omterm{def:b2:funcseq:def}{uniformément} \emph{sur
tout $\R$} vers $f$, alors $f$ est un polynôme. \emph{(Pour $m, n$
grands, $P_n - P_m$ est un polynôme borné sur $\R$, donc constant~;
ainsi la suite se stabilise modulo des constantes.)}
\end{exercise}

\begin{exercise}[$\star\star\star$]\label{exo:b2:funcseq:9}
(Une fonction \omterm{def:b2:metric:continuity}{continue}, nulle part dérivable --- guidé) Soit
$\varphi$ la distance à l'entier le plus proche ($1$-périodique,
$\norm{\varphi}_\infty = \frac12$, $1$-lipschitzienne) et
\[
W(x) = \sum_{n=0}^{\infty} \Bigl(\frac{3}{4}\Bigr)^{\!n}
\varphi(4^n x) .
\]
Prouver~: (a) $W$ est \omterm{def:b2:metric:continuity}{continue} sur $\R$ (\omterm{def:b2:funcseq:series}{convergence normale})~; (b)
pour tout $x$ et tout $m$, en choisissant $h_m = \pm\frac12\cdot 4^{-m}$
avec le signe rendant $\varphi$ affine sur le segment de $4^m x$
à $4^m(x + h_m)$, le taux d'accroissement vérifie
\[
\Bigl|\frac{W(x + h_m) - W(x)}{h_m}\Bigr| \geq 3^m -
\sum_{n<m} 3^n \geq \frac{3^m + 1}{2} \xrightarrow[m\to\infty]{}
\infty
\]
(les termes $n > m$ s'annulent par périodicité~; le terme $n = m$
contribue exactement $3^m$~; les termes $n < m$ sont bornés par la
propriété de Lipschitz). Conclure que $W$ n'est dérivable nulle part.
\end{exercise}

\begin{exercise}[$\star$]\label{exo:b2:funcseq:10}
Soit $u_n(x) = (-1)^n x^n(1 - x)$ sur $\intcc{0}{1}$. Montrer que
$\sum u_n$ \omterm{def:b2:integration:improper}{converge} \omterm{def:b2:funcseq:def}{simplement} sur $\intcc{0}{1}$ et calculer sa
somme~; montrer que la convergence est uniforme sur $\intcc{0}{1}$
\emph{(majorer le reste $R_N(x) = \sum_{n > N} u_n(x)$, une
queue géométrique, par son premier terme et maximiser $x^{N+1}(1-x)$)}
mais \emph{non} normale \emph{(calculer $\norm{u_n}_\infty$)}~:
la \omterm{def:b2:funcseq:def}{convergence uniforme} est strictement plus faible que la \omterm{def:b2:funcseq:series}{convergence
normale}. Contraster avec $\sum x^n(1-x)$, dont la somme est
discontinue en $1$~: là, même l'uniformité échoue.
\end{exercise}

\begin{exercise}[$\star\star$]\label{exo:b2:funcseq:11}
Soit $f_n \to f$ \omterm{def:b2:funcseq:def}{uniformément} sur un \omterm{def:b2:metric:def}{espace métrique} $X$, chaque $f_n$
\omterm{def:b2:metric:continuity}{continue}, et soit $x_n \to x$ dans $X$. Prouver $f_n(x_n) \to
f(x)$. Montrer par un exemple sur $X = \intcc{0}{1}$ que la
\omterm{def:b2:funcseq:def}{convergence simple} ne suffit pas, même avec $f$ \omterm{def:b2:metric:continuity}{continue}
\emph{(utiliser les bosses $g_n$ de l'\cref{exo:b2:funcseq:1} et $x_n =
\frac{1}{\sqrt{2n}}$)}.
\end{exercise}

\begin{exercise}[$\star\star\star$]\label{exo:b2:funcseq:12}
(Une équation intégrale de Volterra par séries) Pour $f \in
C(\intcc{0}{1})$ définir $Tf(x) = \int_0^x f(t)\,\dd t$.
\begin{enumerate}
  \item Montrer par récurrence que pour $n \geq 1$~:
        \[
        T^n f(x) = \int_0^x \frac{(x -
        t)^{n-1}}{(n-1)!}\,f(t)\,\dd t,
        \qquad
        \norm{T^n f}_\infty \leq \frac{\norm f_\infty}{n!} .
        \]
  \item En déduire que $S = \sum_{n\geq0} T^n f$ \omterm{def:b2:integration:improper}{converge} \omterm{def:b2:funcseq:series}{normalement}
        sur $\intcc{0}{1}$ et résout l'équation intégrale $S = f
        + TS$.
  \item Vérifier que $S(x) = f(x) + \int_0^x \eu^{x-t}f(t)\,\dd
        t$ résout la même équation, et prouver l'unicité des
        solutions \omterm{def:b2:metric:continuity}{continues} \emph{(si $S = TS$ alors
        $\norm{S}_\infty \leq \norm{T^nS}_\infty \to 0$)}~:
        conclure la forme close de la somme.
\end{enumerate}
\end{exercise}

\section{Problème~: vitesses d'approximation et théorème de Korovkin}

\begin{problem}\label{pb:b2:funcseq:1}
La preuve de Bernstein du \cref{thm:b2:funcseq:weierstrass} recèle
deux trésors. D'abord, elle est \emph{quantitative}~: la vitesse à
laquelle $B_nf \to f$ est régie par le module de continuité de $f$,
la vitesse optimale étant atteinte par $\abs{x - \frac12}$. Ensuite,
elle est \emph{structurelle}~: tout ce qui a compté est que $B_n$ est
un opérateur linéaire positif se comportant bien sur $1$, $x$, $x^2$
--- cette observation, isolée, est le \emph{théorème de
Korovkin}\index{théorème de Korovkin}. Ce problème démontre les deux,
et se clôt sur l'asymptotique exacte de Voronovskaya. Partout,
$f \in C(\intcc{0}{1})$, $M = \norm f_\infty$, $p_k(x) = \binom
nk x^k(1-x)^{n-k}$, et $e_j$ désigne $x \mapsto x^j$.

\medskip
\noindent\textbf{Partie~I --- L'opérateur de Bernstein.}
\begin{enumerate}
  \item Montrer que $B_n$ est linéaire, \emph{positif} ($f \geq 0
        \Rightarrow B_nf \geq 0$), donc monotone ($f \leq g
        \Rightarrow B_nf \leq B_ng$), avec
        $\norm{B_nf}_\infty \leq \norm f_\infty$, et que
        $B_nf$ interpole $f$ aux deux extrémités.
  \item Redémontrer les identités $B_n e_0 = e_0$, $B_n e_1 =
        e_1$ et $B_n e_2 = e_2 + \frac{e_1 - e_2}{n}$
        \emph{(dériver $(x + y)^n$ deux fois et poser $y = 1 -
        x$)}.
  \item En déduire l'identité de variance $\sum_k \bigl(\frac kn -
        x\bigr)^2 p_k(x) = \frac{x(1-x)}{n}$ et, par
        Cauchy--Schwarz, la borne de premier moment
        \[
        \sum_{k=0}^{n}\Bigl|\frac kn - x\Bigr|\,p_k(x)
        \leq \sqrt{\frac{x(1-x)}{n}} \leq \frac{1}{2\sqrt n} .
        \]
  \item Montrer que si $f$ est convexe alors $B_nf \geq f$ sur
        $\intcc{0}{1}$ \emph{(inégalité de Jensen finie pour les
        poids $p_k(x)$)}.
  \item (Borne de comptage de Tchebychev, reformulée) Pour $\delta > 0$
        montrer
        \[
        \sum_{\abs{k/n - x} > \delta} p_k(x)
        \leq \frac{x(1-x)}{n\delta^2}
        \leq \frac{1}{4n\delta^2} ,
        \]
        et donner la lecture probabiliste~: $B_nf(x)$ moyenne
        $f$ sur une moyenne empirique binomiale qui se concentre en $x$.
\end{enumerate}

\medskip
\noindent\textbf{Partie~II --- Vitesses~: le module de continuité.}
Pour $\delta > 0$ posons $\omega(\delta) = \sup\{\abs{f(s) - f(t)}
: s, t \in \intcc{0}{1},\ \abs{s - t} \leq \delta\}$.
\begin{enumerate}[resume]
  \item Montrer~: $\omega$ est finie, croissante,
        $\omega(\delta) \to 0$ quand $\delta \to 0^+$ (Heine),
        sous-additive ($\omega(\delta_1 + \delta_2) \leq
        \omega(\delta_1) + \omega(\delta_2)$), et
        $\omega(\lambda\delta) \leq (1 +
        \lambda)\,\omega(\delta)$ pour tout $\lambda > 0$.
  \item Prouver l'estimation maîtresse, pour tout $\delta > 0$~:
        \[
        \abs{B_nf(x) - f(x)}
        \leq \sum_k \omega\Bigl(\Bigl|\frac kn -
        x\Bigr|\Bigr)p_k(x)
        \leq \Bigl(1 + \frac1\delta\sum_k\Bigl|\frac kn -
        x\Bigr|p_k(x)\Bigr)\,\omega(\delta) .
        \]
  \item Choisir $\delta = n^{-1/2}$ et conclure le
        \emph{théorème de Weierstrass quantitatif}~:
        \[
        \norm{B_nf - f}_\infty \leq
        \frac32\,\omega\Bigl(\frac{1}{\sqrt n}\Bigr)
        \xrightarrow[n\to\infty]{} 0 .
        \]
  \item En déduire les vitesses~: $\norm{B_nf - f}_\infty \leq
        \frac{3L}{2\sqrt n}$ pour $f$ $L$-lipschitzienne, et $\leq
        \frac32 C n^{-\alpha/2}$ pour $f$ $\alpha$-höldérienne
        ($\abs{f(s) - f(t)} \leq C\abs{s-t}^\alpha$).
  \item (L'exemple optimal --- une identité binomiale) Pour $m
        \geq 1$ prouver
        \[
        \sum_{k=m+1}^{2m} (k - m)\binom{2m}{k}
        = \frac{m}{2}\binom{2m}{m},
        \qquad\text{d'où}\qquad
        \sum_{k=0}^{2m}\abs{k - m}\binom{2m}{k}
        = m\binom{2m}{m}
        \]
        \emph{(utiliser $k\binom{2m}k = 2m\binom{2m-1}{k-1}$ et la
        symétrie de la ligne binomiale, qui donne
        $\sum_{j=m}^{2m-1}\binom{2m-1}{j} = 2^{2m-2}$)}.
  \item Pour $f(t) = \abs{t - \frac12}$ en déduire la valeur exacte
        et son asymptotique (binomial central,
        \cref{ex:b2:comparison:centralbinomial})~:
        \[
        B_{2m}f\Bigl(\frac12\Bigr) - f\Bigl(\frac12\Bigr)
        = \frac{\binom{2m}{m}}{2\cdot4^{m}}
        \;\sim\; \frac{1}{2\sqrt{\pi m}} :
        \]
        la vitesse $\omega(n^{-1/2})$ de la question~8 est atteinte
        (à une constante près) --- pour $f$ seulement \omterm{def:b2:metric:continuity}{continue},
        le $n^{-1/2}$ de Bernstein est honnête.
\end{enumerate}

\medskip
\noindent\textbf{Partie~III --- Théorème de Korovkin.} Soit $(L_n)$
une suite d'opérateurs \emph{linéaires positifs} de
$C(\intcc{0}{1})$ dans lui-même telle que $L_ne_j \to e_j$
\omterm{def:b2:funcseq:def}{uniformément} pour $j = 0, 1, 2$.
\begin{enumerate}[resume]
  \item Montrer qu'un opérateur linéaire positif $L$ est monotone et
        vérifie $\abs{Lf} \leq L\abs f$ ponctuellement.
  \item Montrer~: pour tout $\varepsilon > 0$ il existe $\delta > 0$
        tel que pour \emph{tous} $s, x \in \intcc{0}{1}$~:
        \[
        \abs{f(s) - f(x)} \leq \varepsilon +
        \frac{2M}{\delta^2}(s - x)^2
        \]
        \emph{(traiter $\abs{s - x} \leq \delta$ par Heine et
        $\abs{s-x} > \delta$ par la borne grossière $2M$)}.
  \item Fixer $x$, appliquer $L_n$ à l'inégalité de la question~13
        dans la variable $s$, et en déduire
        \[
        \abs{L_nf(x) - f(x)\,L_ne_0(x)}
        \leq \varepsilon\,L_ne_0(x) + \frac{2M}{\delta^2}
        \bigl(L_ne_2(x) - 2x\,L_ne_1(x) + x^2 L_ne_0(x)\bigr).
        \]
  \item Montrer que $\sup_x \bigl(L_ne_2(x) - 2x\,L_ne_1(x) + x^2
        L_ne_0(x)\bigr) \to 0$, puis assembler le
        \emph{théorème de Korovkin}~: $L_nf \to f$ \omterm{def:b2:funcseq:def}{uniformément} pour
        \emph{toute} $f \in C(\intcc{0}{1})$.
  \item Vérifier que $(B_n)$ satisfait les hypothèses de Korovkin~:
        Weierstrass une troisième fois, à partir de trois monômes.
  \item Soit $I_n$ l'opérateur d'interpolation affine par morceaux
        aux nœuds $\frac kn$. Montrer que $I_n$ est
        linéaire positif, $I_ne_0 = e_0$, $I_ne_1 = e_1$, et
        $\norm{I_ne_2 - e_2}_\infty = \frac{1}{4n^2}$
        \emph{(sur chaque maille l'erreur de l'interpolation affine
        de $t^2$ est $(t - a)(b - t)$)}. Conclure par Korovkin~: les
        interpolants polygonaux convergent \omterm{def:b2:funcseq:def}{uniformément} pour toute
        fonction $f$ \omterm{def:b2:metric:continuity}{continue}.
\end{enumerate}

\medskip
\noindent\textbf{Partie~IV --- Dividendes~: densité, moments,
dérivées.}
\begin{enumerate}[resume]
  \item Montrer que les polynômes à coefficients \emph{rationnels}
        sont denses dans $\bigl(C(\intcc{0}{1}),
        \norm\cdot_\infty\bigr)$~: cet \omterm{def:b2:nvs:banach}{espace de Banach} est
        séparable.
  \item (Les moments déterminent la fonction) Soit $f \in
        C(\intcc{0}{1})$ avec $\int_0^1 f(t)\,t^n \dd t = 0$
        pour tout $n \in \N$. Montrer $\int_0^1 f P = 0$ pour tout
        polynôme, puis $\int_0^1 f^2 = 0$, puis $f = 0$.
  \item Prouver l'identité de dérivée
        \[
        (B_nf)'(x) = n\sum_{k=0}^{n-1}\Bigl(
        f\Bigl(\frac{k+1}{n}\Bigr) -
        f\Bigl(\frac kn\Bigr)\Bigr)\,
        \binom{n-1}{k}x^k(1-x)^{n-1-k}
        \]
        \emph{(dériver $p_k$ et réindexer --- une \omterm{thm:b2:series:abel}{transformation
        d'Abel})}.
  \item Supposons $f$ de classe $C^1$. En utilisant le théorème des
        accroissements finis dans chaque incrément et en comparant
        avec $B_{n-1}(f')$, montrer $(B_nf)' \to f'$ \omterm{def:b2:funcseq:def}{uniformément} sur
        $\intcc{0}{1}$. En déduire~: pour $f \in C^1$ il existe des
        polynômes convergeant vers $f$ \emph{en même temps que} leurs
        dérivées.
  \item Supposons $f$ de classe $C^2$. Par Taylor--Lagrange en $x$
        montrer
        \[
        \abs{B_nf(x) - f(x)} \leq
        \frac{\norm{f''}_\infty}{2}\cdot\frac{x(1-x)}{n}
        \leq \frac{\norm{f''}_\infty}{8n} :
        \]
        la régularité améliore la vitesse de $n^{-1/2}$ à
        $n^{-1}$.
\end{enumerate}

\medskip
\noindent\textbf{Partie~V --- Saturation~: le théorème de
Voronovskaya.}
\begin{enumerate}[resume]
  \item Prouver l'identité de quatrième moment
        \[
        \sum_k (k - nx)^4 p_k(x)
        = nx(1-x)\bigl(1 + 3(n-2)x(1-x)\bigr) \leq n^2
        \quad (n \geq 1)
        \]
        \emph{(développer $k^4$ en factorielles décroissantes
        $k(k-1)\cdots$ et utiliser deux fois de plus l'astuce de
        dérivation de la question~2)}.
  \item (Voronovskaya) Soit $f$ de classe $C^2$ et $x \in
        \intcc{0}{1}$. En écrivant $f(t) = f(x) + f'(x)(t-x) +
        \frac{f''(x)}2(t-x)^2 + \eta(t)(t-x)^2$ avec $\eta$
        bornée et $\eta(t) \to 0$ quand $t \to x$, prouver
        \[
        n\bigl(B_nf(x) - f(x)\bigr)
        \xrightarrow[n\to\infty]{}
        \frac{x(1-x)}{2}\,f''(x)
        \]
        \emph{(scinder la somme des $\eta$ en $\abs{t - x} \leq
        \delta$~; contrôler la partie lointaine avec la question~23)}.
        Ainsi l'erreur de la question~22 est exacte en \omterm{def:b2:structures:generated}{ordre}
        \emph{et} en constante~: $B_n$ \emph{sature} à $\frac1n$, si
        régulière que soit $f$ --- comparer
        l'\cref{exo:b2:funcseq:7}.
  \item Synthèse. En une phrase chacun~: (i) ce que la positivité
        seule a apporté (Parties~I et III)~; (ii) où la compacité
        de $\intcc{0}{1}$ est intervenue dans chaque partie~; (iii)
        pourquoi trois fonctions test suffisent dans le théorème de
        Korovkin~; (iv) le compromis que fait Bernstein (un
        $n^{-1/2}$ robuste pour $f$ irrégulière, mais un plafond en
        $\frac1n$ pour $f$ régulière), et quel chapitre de ce livre
        jouera le même jeu avec des polynômes trigonométriques.
\end{enumerate}
\end{problem}
